\documentclass[11pt]{article}
\usepackage[margin=1in]{geometry}
\usepackage{amsmath}
\usepackage{amsfonts}
\usepackage{amssymb}
\usepackage{algorithm}
\usepackage{algorithmic}
\usepackage{graphicx}
\usepackage{cite}
\usepackage{url}
\usepackage{booktabs}
\usepackage{array}
\usepackage{xepersian}
\settextfont{XB Niloofar}

% Define Persian commands
\newcommand{\فصل}[1]{\section{#1}}
\newcommand{\قسمت}[1]{\subsection{#1}}

\title{حافظه‌نهانی هوشمند برای شبکه‌های وسایل نقلیه نرم‌افزار محور: رویکردی مبتنی بر یادگیری فدرال و یادگیری تقویتی}
\author{گزارش پیشرفت رساله دکتری}
\date{\today}

\begin{document}

\maketitle

\فصل{پیشنهاد پژوهشی رساله‌ی دکتری}

\قسمت{تعریف مسئله}

شبکه‌های خودروی موردی (VANETs) به عنوان زیرساخت حیاتی برای فعال‌سازی سیستم‌های حمل‌ونقل هوشمند، رانندگی خودکار و کاربردهای خودروی بلادرنگ ظهور کرده‌اند. ادغام اصول شبکه‌سازی نرم‌افزار محور (SDN) با VANETs کنترل متمرکز و قابلیت برنامه‌ریزی را فراهم می‌کند که مدیریت کارآمد منابع و زمان‌بندی وظایف را امکان‌پذیر می‌سازد.

با این حال، ماهیت پویای محیط‌های خودروی که با تحرک بالا، اتصال متناوب و الزامات تأخیر سخت مشخص می‌شود، چالش‌های قابل توجهی برای حفظ تضمین‌های کیفیت سرویس (QoS) ایجاد می‌کند.

این تحقیق مسئله اساسی حافظه‌نهانی هوشمند در VANETs فعال‌شده با SDN را برای بهینه‌سازی عملکرد زمان‌بندی وظایف در حین رعایت محدودیت‌های ضرب‌الاجل سخت مورد بررسی قرار می‌دهد.

ما معماری دو مرحله‌ای نوآورانه‌ای پیشنهاد می‌کنیم که یادگیری فدرال (FL) را برای تصمیمات پیش‌قرارگیری حافظه نهان با یادگیری تقویتی (RL) برای سیاست‌های پویای اخراج حافظه نهان ترکیب می‌کند.

هدف اصلی حداکثرسازی نرخ برخورد حافظه نهان (هدف $\geq 74\%$) در حین حفظ نرخ موفقیت بالای وظایف تحت محدودیت‌های بلادرنگ است.

با توجه به VANET متشکل از $N$ خودرو و $M$ واحد کنار جاده (RSU) که به عنوان گره‌های محاسبات مه عمل می‌کنند، سیستم زمان گسسته‌ای را در نظر می‌گیریم که خودروها وظایف محاسباتی با ضرب‌الاجل‌های سخت تولید می‌کنند.

هر وظیفه $\tau_i$ با مشخصات زیر تعریف می‌شود:

\begin{equation}
\tau_i = \langle id_i, type_i, priority_i, size_i, deadline_i, origin_{RSU}, destination_{RSU} \rangle
\end{equation}

چالش طراحی سیستم حافظه‌نهانی هوشمند است که سه هدف عملکردی حیاتی را مورد توجه قرار دهد.

اول، سیستم باید نرخ برخورد حافظه نهان را که به صورت $\eta = \frac{\text{cache hits}}{\text{total requests}}$ تعریف می‌شود، با آستانه هدف حداقل 74\% حداکثر کند. این معیار مستقیماً بر کارایی سیستم تأثیر می‌گذارد و با حذف نیاز به محاسبه مجدد وظایف پردازش‌شده قبلی، افزونگی محاسباتی و سربار شبکه را کاهش می‌دهد.

دوم، سیستم باید نرخ موفقیت بالای وظایف $\sigma = \frac{\text{tasks meeting deadlines}}{\text{total tasks}}$ حداقل 80\% را حفظ کند تا اطمینان حاصل شود که اکثریت وظایف خودروی علی‌رغم شرایط پویای شبکه در ضرب‌الاجل‌های مشخص‌شده تکمیل می‌شوند.

سوم، سیستم باید میانگین زمان تکمیل $\bar{T} = \frac{1}{|\mathcal{T}|} \sum_{i \in \mathcal{T}} T_i$ را در تمام وظایف تکمیل‌شده کمینه کند و بدین ترتیب پاسخگویی کلی سیستم و تجربه کاربر در کاربردهای خودروی بلادرنگ را بهبود بخشد.

\قسمت{گام‌های تحقیقاتی}

این پژوهش در چندین مرحله تعریف‌شده انجام می‌پذیرد که هر مرحله اهداف مشخص و قابل اندازه‌گیری دارد.

\textbf{گام اول - تحلیل و طراحی معماری سیستم:} در این مرحله، معماری جامع سیستم حافظه‌نهانی هوشمند طراحی می‌شود. این معماری شامل چهار لایه اصلی است که عبارتند از لایه شبکه فیزیکی (شامل خودروها و RSUها)، لایه کنترل SDN (شامل کنترلر متمرکز و الگوریتم‌های زمان‌بندی)، لایه یادگیری هوشمند (شامل سیستم‌های FL و RL) و لایه مدیریت داده (شامل سیستم‌های حافظه نهان محلی). در این مرحله همچنین تعامل بین اجزای مختلف سیستم، پروتکل‌های ارتباطی و جریان داده‌ها تعریف می‌شود.

\textbf{گام دوم - مدل‌سازی ریاضی و فرمول‌بندی مسائل:} این مرحله شامل توسعه مدل‌های ریاضی دقیق برای تمامی اجزای سیستم است. مسئله زمان‌بندی وظایف به صورت برنامه‌ریزی خطی عدد صحیح مختلط (MILP) فرمول‌بندی می‌شود. الگوریتم یادگیری فدرال با استفاده از FedAvg و شبکه‌های عصبی چندلایه مدل‌سازی می‌گردد. سیستم یادگیری تقویتی نیز به صورت فرآیند تصمیم‌گیری مارکوف (MDP) با فضای حالت 9 بعدی و 5 عمل ممکن تعریف می‌شود.

\textbf{گام سوم - پیاده‌سازی اجزای سیستم:} در این مرحله، تمامی اجزای طراحی‌شده در مراحل قبل پیاده‌سازی می‌شوند. این شامل توسعه مدیر حافظه نهان توزیع‌شده با قابلیت‌های LFU/LRU، پیاده‌سازی مشاوران FL محلی در هر RSU، توسعه عامل RL سراسری با معماری DQN، ایجاد سیستم پیش‌بینی مسیر خودروها با XGBoost و پیاده‌سازی الگوریتم‌های زمان‌بندی MILP است.

\textbf{گام چهارم - ارزیابی جامع و آزمایش‌های تطبیقی:} این مرحله شامل طراحی و اجرای آزمایش‌های جامع برای ارزیابی عملکرد سیستم است. آزمایش‌ها با استفاده از داده‌های مصنوعی واقع‌گرایانه انجام می‌پذیرد و شامل مقایسه با روش‌های پایه مختلف نظیر LFU، LRU، اخراج تصادفی و یادگیری فدرال متمرکز است. معیارهای ارزیابی شامل نرخ برخورد حافظه نهان، نرخ موفقیت وظایف، میانگین زمان تکمیل و سربار محاسباتی است.

\textbf{گام پنجم - بهینه‌سازی و تنظیم پارامترها:} در این مرحله نهایی، پارامترهای مختلف سیستم شامل ضریب زوال محبوبیت، فرکانس تجمیع FL، اندازه حافظه نهان، پارامترهای شبکه عصبی و تنظیمات RL بهینه‌سازی می‌شوند. همچنین تحلیل حساسیت برای شناسایی پارامترهای حیاتی انجام می‌پذیرد و راهکارهای بهبود عملکرد ارائه می‌گردد.

\فصل{فعالیت‌های انجام شده تاکنون}

\قسمت{معماری کلی سیستم}

سیستم حافظه‌نهانی هوشمند پیشنهادی بر اساس معماری چهارلایه‌ای پیشرفته طراحی شده که ترکیبی منحصربه‌فرد از تکنولوژی‌های شبکه‌سازی نرم‌افزار محور، یادگیری فدرال و یادگیری تقویتی را ارائه می‌دهد. این معماری به گونه‌ای طراحی شده که قابلیت مقیاس‌پذیری، انطباق‌پذیری و کارایی بالا را در محیط‌های پویای خودروی تضمین کند.

\begin{figure}[h]
\centering
\includegraphics[width=0.8\textwidth]{system_architecture.png}
\caption{معماری کلی سیستم حافظه‌نهانی هوشمند}
\end{figure}

\textbf{لایه اول - زیرساخت شبکه فیزیکی:}

این لایه شامل اجزای فیزیکی شبکه است که پایه و اساس کل سیستم را تشکیل می‌دهد:

\textit{واحدهای کنار جاده (RSUs):} شبکه شامل 5 واحد RSU استراتژیک قرار گرفته در نقاط کلیدی شهری است. هر RSU به عنوان یک گره محاسباتی مه (Fog Node) عمل کرده و دارای مشخصات زیر است:
\begin{itemize}
\item منابع محاسباتی: 8 هسته CPU با فرکانس 2.4 GHz
\item حافظه RAM: 16 GB DDR4
\item فضای ذخیره‌سازی: 1 TB SSD با سرعت خواندن/نوشتن بالا
\item حافظه نهان محلی: ظرفیت 100 ورودی با پشتیبانی از الگوریتم‌های LRU/LFU
\item رادیو DSRC/5G: برد ارتباطی 300 متر با پهنای باند 20 MHz
\item اتصال شبکه: Ethernet Gigabit به کنترلر مرکزی
\end{itemize}

\textit{خودروهای هوشمند:} شبکه شامل 50 خودرو متحرک است که هر کدام دارای ویژگی‌های زیر هستند:
\begin{itemize}
\item واحد ارتباطی خودرو (OBU): پشتیبانی از DSRC و 5G-V2X
\item سیستم موقعیت‌یابی: GPS با دقت زیر 1 متر
\item الگوی تحرک: سرعت متغیر 0-120 km/h بر اساس مدل شهری واقعی
\item تولید وظایف: فرآیند پواسون با نرخ λ = 0.3 وظیفه در ثانیه
\item انواع وظایف: پردازش تصویر، تشخیص اشیاء، مسیریابی، تحلیل ترافیک
\end{itemize}

\textbf{لایه دوم - کنترل و هماهنگی SDN:}

این لایه قلب سیستم محسوب شده و مسئولیت کنترل متمرکز و هماهنگی کلی عملیات را بر عهده دارد:

\textit{کنترلر SDN مرکزی:} واحد کنترل اصلی که بر روی سرور قدرتمند اجرا شده و شامل ماژول‌های زیر است:
\begin{itemize}
\item \textbf{زمان‌بند MILP پیشرفته:} حل مسائل بهینه‌سازی پیچیده با استفاده از حل‌کننده Gurobi
\item \textbf{مدیر حافظه نهان توزیع‌شده:} هماهنگی و همگام‌سازی حافظه‌های نهان RSUها
\item \textbf{ماژول پیش‌بینی مسیر:} الگوریتم XGBoost برای پیش‌بینی مسیر خودروها
\item \textbf{سیستم نظارت بلادرنگ:} جمع‌آوری و تحلیل آمار عملکرد هر 100 میلی‌ثانیه
\item \textbf{واحد تصمیم‌گیری هوشمند:} ادغام نتایج FL و RL برای تصمیمات بهینه
\end{itemize}

\textbf{لایه سوم - یادگیری هوشمند و تطبیقی:}

این لایه مغز هوشی سیستم محسوب شده و شامل دو زیرسیستم یادگیری پیشرفته است:

\textit{زیرسیستم یادگیری فدرال توزیع‌شده:}
\begin{itemize}
\item \textbf{مشاوران FL محلی:} هر RSU یک مشاور FL مستقل با شبکه عصبی MLP اجرا می‌کند
\item \textbf{استخراج ویژگی پیشرفته:} بردار 12 بعدی شامل فراوانی، اولویت، اندازه و متادیتای زمانی
\item \textbf{آموزش محلی حفظ حریم خصوصی:} SGD محلی بدون انتقال داده خام
\item \textbf{تجمیع FedAvg:} الگوریتم میانگین‌گیری وزنی هر 10 گام زمانی
\item \textbf{پیش‌قرارگیری هوشمند:} انتخاب top-K وظایف محبوب بر اساس پیش‌بینی مدل
\end{itemize}

\textit{زیرسیستم یادگیری تقویتی سراسری:}
\begin{itemize}
\item \textbf{عامل DQN مرکزی:} شبکه عصبی عمیق با معماری 9-64-64-5
\item \textbf{فضای حالت جامع:} 9 ویژگی تجمیع‌شده از کل سیستم
\item \textbf{فضای عمل انطباقی:} 5 استراتژی اخراج مختلف
\item \textbf{تابع پاداش چندهدفه:} متعادل‌کننده نرخ برخورد، موفقیت و زمان پاسخ
\item \textbf{بازپخت تجربه:} حافظه 10000 تجربه برای آموزش پایدار
\end{itemize}

\textbf{لایه چهارم - مدیریت داده و حافظه نهان:}

این لایه مسئول ذخیره‌سازی، مدیریت و بازیابی کارآمد داده‌ها است:

\textit{سیستم حافظه نهان توزیع‌شده:}
\begin{itemize}
\item \textbf{ساختار داده بهینه:} Hash Table با O(1) دسترسی و LRU/LFU tracking
\item \textbf{مدیریت TTL پیشرفته:} زمان انقضای تطبیقی بر اساس الگوهای دسترسی
\item \textbf{سیاست‌های اخراج هیبریدی:} ترکیب LRU، LFU و تصادفی با وزن‌دهی پویا
\item \textbf{همگام‌سازی توزیع‌شده:} پروتکل consistency برای حفظ یکپارچگی داده‌ها
\item \textbf{آمارگیری بلادرنگ:} ردیابی hit rate، access frequency و cache utilization
\end{itemize}

\textit{سیستم جستجوی توزیع‌شده:}
\begin{itemize}
\item \textbf{جستجوی چندمرحله‌ای:} محلی → مجاور → سراسری
\item \textbf{تکرار خودکار:} کپی‌برداری هوشمند از محتوای محبوب
\item \textbf{بالانس بار:} توزیع یکنواخت درخواست‌ها بین RSUها
\item \textbf{مدیریت شکست:} مکانیزم fallback و recovery خودکار
\end{itemize}

\قسمت{معماری یادگیری فدرال}

سیستم یادگیری فدرال طراحی‌شده دارای معماری کاملاً توزیع‌شده است که حریم خصوصی داده‌ها را حفظ می‌کند:

\textbf{کلاینت‌های محلی (مشاوران FL):}
هر RSU یک مشاور FL مستقل را اجرا می‌کند که شامل:
- شبکه عصبی MLP با 12 ورودی و 2 لایه مخفی
- ماژول استخراج ویژگی از وظایف محلی
- سیستم آموزش محلی با SGD
- مکانیزم ذخیره و بازیابی وزن‌های مدل

بردار ویژگی 12 بعدی شامل موارد زیر است:
- فراوانی نرمال‌شده درخواست: $freq_i / freq_{max}$
- فراوانی لگاریتمی: $\log_{10}(1 + freq_i) / \log_{10}(1 + freq_{max})$
- اولویت نرمال‌شده: $priority_i / 3$
- اندازه داده نرمال‌شده: $size_i / 1000$
- پرچم محلی‌بودن: $locality_i$
- سرعت خودرو نرمال‌شده: $speed_i / 120$
- بار RSU: $load_i$
- ضریب زوال محبوبیت: $decay_i$
- کدگذاری یک-گرم نوع وظیفه (4 بعد)

فرمول‌بندی ریاضی برای به‌روزرسانی‌های مدل محلی:

\begin{equation}
\theta_r^{(t+1)} = \theta_r^{(t)} - \eta \nabla_{\theta_r} \mathcal{L}_r(\theta_r^{(t)}, \mathcal{D}_r^{(t)})
\end{equation}

تابع زیان برای آموزش محلی:

\begin{equation}
\mathcal{L}_r(\theta_r) = \frac{1}{|\mathcal{B}_r|} \sum_{i \in \mathcal{B}_r} \left[ y_i \log(f_{\theta_r}(\mathbf{x}_i)) + (1-y_i) \log(1-f_{\theta_r}(\mathbf{x}_i)) \right]
\end{equation}

\textbf{سرور تجمیع مرکزی:}
کنترلر SDN نقش سرور FL را ایفا می‌کند و شامل:
- الگوریتم FedAvg برای تجمیع وزن‌های مدل
- سیستم هماهنگی دورهای (هر 10 گام زمانی)
- مکانیزم پخش وزن‌های سراسری به تمام RSUها
- نظارت بر همگرایی مدل‌ها

تجمیع مدل سراسری الگوریتم FedAvg را دنبال می‌کند:

\begin{equation}
\theta^{(t+1)} = \sum_{r=1}^{M} \frac{|\mathcal{D}_r|}{|\mathcal{D}|} \theta_r^{(t+1)}
\end{equation}

\textbf{فرآیند پیش‌قرارگیری:}
پس از آموزش و تجمیع، هر مشاور FL لیستی از کلیدهای محبوب را تولید می‌کند:
- محاسبه امتیاز محبوبیت برای هر وظیفه
- انتخاب top-K وظایف با بالاترین امتیاز
- پیش‌قرارگیری در تمام RSUها برای دسترسی سریع

\قسمت{معماری یادگیری تقویتی}

عامل یادگیری تقویتی سراسری با معماری DQN طراحی شده که شامل:

\textbf{فضای حالت 9 بعدی:}
عامل RL وضعیت تجمیع‌شده سیستم را مشاهده می‌کند:

\begin{equation}
s_t = \langle u_t, h_t, \bar{T}_t, \bar{p}_t, \bar{s}_t, l_t, \bar{v}_t, \bar{\rho}_t, \bar{\delta}_t \rangle
\end{equation}

که در آن:
- $u_t$: نسبت استفاده از حافظه نهان
- $h_t$: نرخ برخورد فعلی
- $\bar{T}_t$: میانگین زمان تکمیل
- $\bar{p}_t$: میانگین اولویت وظایف
- $\bar{s}_t$: میانگین اندازه وظایف
- $l_t$: نسبت محلی‌بودن
- $\bar{v}_t$: میانگین سرعت خودروها
- $\bar{\rho}_t$: میانگین بار RSUها
- $\bar{\delta}_t$: میانگین زوال محبوبیت

\textbf{فضای عمل:}
عامل می‌تواند یکی از 5 عمل زیر را انتخاب کند:
- عدم اقدام (noop)
- اخراج کمترین استفاده اخیر (evict\_lru)
- اخراج تصادفی (evict\_random)  
- اخراج کمترین فراوانی (evict\_lfu)
- کاهش TTL (shrink\_ttl)

\textbf{تابع پاداش:}
تابع پاداش چندهدفه طراحی شده که سه معیار عملکرد را متعادل می‌کند:

\begin{equation}
R(s_t, a_t) = \alpha \cdot \Delta h_t + \beta \cdot \Delta \sigma_t - \gamma \cdot \Delta \bar{T}_t
\end{equation}

\قسمت{مدل‌سازی ریاضی}

\textbf{بهینه‌سازی زمان‌بندی وظایف:}
مسئله زمان‌بندی وظایف به عنوان برنامه خطی عدد صحیح مختلط (MILP) فرمول‌بندی شده:

\begin{align}
\text{maximize} \quad & \sum_{i \in \mathcal{T}} x_i \cdot w_i \\
\text{subject to} \quad & \sum_{i \in \mathcal{T}} x_i \cdot c_i \leq C_{total} \\
& T_i \leq D_i, \quad \forall i \in \mathcal{T} \\
& x_i \in \{0, 1\}, \quad \forall i \in \mathcal{T}
\end{align}

\textbf{مدل یادگیری فدرال:}
هر RSU شبکه عصبی چندلایه (MLP) را نگهداری می‌کند:

\begin{equation}
f_{\theta_r}(\mathbf{x}_i) = \sigma(\mathbf{W}_2 \sigma(\mathbf{W}_1 \mathbf{x}_i + \mathbf{b}_1) + \mathbf{b}_2)
\end{equation}

بردار ویژگی $\mathbf{x}_i \in \mathbb{R}^{12}$ شامل مشخصات متعدد وظیفه و زمینه است.

\قسمت{نتایج آزمایشی}

نتایج آزمایشی نشان‌دهنده عملکرد فوق‌العاده سیستم پیشنهادی است:

\begin{table}[h]
\centering
\begin{tabular}{@{}lccc@{}}
\toprule
روش & نرخ برخورد (\%) & نرخ موفقیت (\%) & زمان اجرا (ثانیه) \\
\midrule
بدون حافظه نهان & 0.0 & 45.2 & 8.2 \\
فقط LFU & 52.3 & 78.1 & 4.1 \\
فقط LRU & 48.7 & 76.8 & 4.3 \\
اخراج تصادفی & 41.2 & 74.5 & 3.9 \\
FL متمرکز & 73.3 & 81.2 & 33.2 \\
\textbf{FL توزیع‌شده + RL} & \textbf{96.9} & \textbf{81.6} & \textbf{5.3} \\
\bottomrule
\end{tabular}
\caption{مقایسه عملکرد روش‌های مختلف}
\end{table}

\فصل{فعالیت‌های پیش‌رو}

\قسمت{تحلیل نظری}

تحقیقات آینده بر توسعه مبانی نظری دقیق متمرکز خواهد بود:
- تضمین‌های همگرایی برای الگوریتم یادگیری فدرال
- کران‌های پشیمانی برای سیاست یادگیری تقویتی
- تحلیل پیچیدگی سیستم مدیریت حافظه نهان توزیع‌شده

\قسمت{توسعه‌های سیستم}

توسعه‌های برنامه‌ریزی‌شده شامل:
- ادغام با ردیابی‌های تحرک واقعی از شبیه‌سازهای SUMO و VANET
- چارچوب‌های بهینه‌سازی چندهدفه که مصرف انرژی را در نظر می‌گیرند
- مکانیزم‌های امنیتی و حریم خصوصی پیشرفته

\قسمت{روش‌های یادگیری پیشرفته}

بررسی روش‌های actor-critic، یادگیری فرا و شبکه‌های عصبی گرافی برای مدل‌سازی الگوهای اتصال پیچیده RSU.

\فصل{جمع‌بندی}

این تحقیق سیستم حافظه‌نهانی هوشمند پیشگامانه‌ای برای VANETs فعال‌شده با SDN ارائه می‌دهد که با موفقیت یادگیری فدرال و یادگیری تقویتی را ترکیب می‌کند تا عملکرد بی‌سابقه‌ای دست یابد.

نتایج آزمایشی ما نرخ برخورد حافظه نهان قابل توجه 96.9\% را نشان می‌دهد که نمایانگر بهبود اساسی نسبت به روش‌های موجود است.

نوآوری‌های کلیدی شامل معماری دو مرحله‌ای نوآورانه، رویکرد یادگیری فدرال حفظ‌کننده حریم خصوصی و چارچوب یادگیری تقویتی سراسری است.

\قسمت{جدول زمانبندی فعالیت‌ها}

\begin{table}[h]
\centering
\begin{tabular}{@{}lcc@{}}
\toprule
فعالیت & زمان شروع & زمان پایان \\
\midrule
طراحی معماری & مهر 1402 & آذر 1402 \\
پیاده‌سازی اولیه & آذر 1402 & بهمن 1402 \\
آزمایش‌های اولیه & بهمن 1402 & اردیبهشت 1403 \\
بهینه‌سازی و تنظیم & اردیبهشت 1403 & مرداد 1403 \\
تحلیل نظری & مرداد 1403 & آبان 1403 \\
نگارش مقالات & آبان 1403 & اسفند 1403 \\
تکمیل رساله & اسفند 1403 & خرداد 1404 \\
\bottomrule
\end{tabular}
\caption{جدول زمان‌بندی فعالیت‌ها}
\end{table}

\قسمت{جدول مقالات پذیرفته‌شده، در حال داوری و در دست تهیه}

\begin{table}[h]
\centering
\begin{tabular}{@{}lcc@{}}
\toprule
عنوان مقاله & مجله/کنفرانس & وضعیت \\
\midrule
Intelligent Caching for VANETs & IEEE TVT & در دست تهیه \\
Federated Learning in Vehicular Networks & Computer Networks & در حال داوری \\
RL-based Cache Management & IEEE INFOCOM & برنامه‌ریزی‌شده \\
\bottomrule
\end{tabular}
\caption{وضعیت مقالات}
\end{table}

\bibliographystyle{IEEEtran}
\bibliography{references}

\end{document}
